\section{Finite polynomial embeddings of symmetric projective bundles}

\begin{lemma}[A graded polynomial presentation of the symmetric algebra]
\label{lem:fc-symmetric-polynomial-presentation}
\leanok
\lean{EnvelopeProperty.SymmetricProjectiveEmbedding.polynomialMap, EnvelopeProperty.SymmetricProjectiveEmbedding.polynomialMap_mem, EnvelopeProperty.SymmetricProjectiveEmbedding.gradedPolynomialMap, EnvelopeProperty.SymmetricProjectiveEmbedding.polynomialMap_surjective, EnvelopeProperty.SymmetricProjectiveEmbedding.gradedPolynomialMap_surjective, EnvelopeProperty.SymmetricProjectiveEmbedding.gradedPolynomialMap_irrelevant_le}
\dcref{rmk:gl-dejong}
Let $M$ be an $R$-module and let $(m_i)_{i\in E}$ generate it.
The homomorphism $R[x_i\mid i\in E]\to\operatorname{Sym}_R M$
sending $x_i$ to the canonical image of $m_i$ preserves every graded
piece and is surjective.
\end{lemma}
\begin{proof}
The degree-one polynomial submodule is the span of the variables, so
its image lies in the image of $M$ in the symmetric algebra.  The
degree-$d$ polynomial submodule is the $d$th power of its degree-one
submodule.  An algebra homomorphism commutes with products and scalar
multiples, hence maps it into the $d$th power of the image of $M$,
which is exactly $\operatorname{Sym}^d_R M$.  For surjectivity, the
subalgebra generated by the images of the $m_i$ contains the image of
all of $M$: its underlying $R$-submodule contains the given generators
and therefore their span.  Every element of the symmetric algebra is
built from scalars and the image of $M$ by finite sums and products,
so the subalgebra is the entire symmetric algebra.
\end{proof}

\begin{theorem}[A finite polynomial embedding of a finite symmetric projective bundle]
\label{thm:fc-symmetric-projective-embedding}
\leanok
\lean{EnvelopeProperty.SymmetricProjectiveEmbedding.embedding, EnvelopeProperty.SymmetricProjectiveEmbedding.embedding_isClosedImmersion, EnvelopeProperty.SymmetricProjectiveEmbedding.embedding_structure, EnvelopeProperty.SymmetricProjectiveEmbedding.exists_embedding, EnvelopeProperty.SymmetricProjectiveEmbedding.isProjectiveFactorization}
\uses{lem:fc-symmetric-polynomial-presentation}
If $M$ is a finite $R$-module, there is a finite-dimensional polynomial
projective space over $R$ and a closed immersion
$\operatorname{Proj}(\operatorname{Sym}_R M)\hookrightarrow\mathbb P^n_R$
whose composite with the original coefficient projection is the
original symmetric-projective-bundle structure morphism.
\end{theorem}
\begin{proof}
Choose a finite generating family for $M$ and append a zero generator
if necessary, so that the indexing set has the form $\{0,\ldots,n\}$.
The preceding lemma gives a surjective graded homomorphism from the
polynomial coordinate ring of $\mathbb P^n_R$ to
$\operatorname{Sym}_R M$.  Surjectivity carries the irrelevant ideal
onto an ideal containing the target irrelevant ideal, so the induced
map on projective spectra is defined.  On every standard homogeneous
affine open, the induced degree-zero localization homomorphism is
surjective: lift a homogeneous numerator, then retain its homogeneous
component of the required degree.  Consequently the projective map
is a closed immersion.  The graded homomorphism is the identity on
$R$; its degree-zero coefficient square therefore proves equality of
the two structure morphisms over $\operatorname{Spec}R$.
\end{proof}

\begin{remark}[The geometric comparison still required for alterations]
\label{rmk:fc-symmetric-embedding-alteration-comparison}
The theorem concerns the actual projective spectrum of a finite
module.  To use it for a projective morphism defined by a finite
quasi-coherent projective bundle, one must identify that whole bundle
over an affine base with the projective spectrum of its global
module.  To obtain an alteration in multiprojective coordinate form,
one must additionally construct its product embedding and identify
the coordinate ideal and all affine coordinate charts.  These are
separate geometric comparisons.
\end{remark}
